\chapter*{Introduction}
\addcontentsline{toc}{chapter}{Introduction}

\section*{Objective and structure of the thesis}
\addcontentsline{toc}{section}{Objective and structure of the thesis}


The objective of this thesis is to quantify how the astrophysical uncertainties 
on the dark matter (DM) density profile of the Milky Way affect the constraints obtained from 
indirect searches towards the Galactic Center. 
The gamma-ray flux expected from DM annihilation depends 
on the density profile through the $J$-factor, a line-of-sight integral of the squared density: as a consequence, different profiles lead to very different constraints.
In order to show this, we compare a set of representative DM density profiles and confront their upper limits on the annihilation cross section with a simple thermal relic model.

The thesis is organized as follows:

\begin{itemize}
    
    \item In the first chapter, we present the main properties of DM, together with its most relevant candidates: WIMPs. Then we discuss the most significant astrophysical evidence, on scales ranging from galaxies and galaxy clusters to the entire Universe. In addition, we briefly introduce theories of modified gravity, developed as an alternative to the DM hypothesis. Finally, we illustrate the theoretical framework for DM indirect searches and describe the main astrophysical environments investigated.
          The primary source for this chapter is \emph{Dark Matter} \cite{cirelli2025darkmatter}, the most complete review on DM written to date.
    
    \item In the second chapter, we focus on the Galactic Center, the most promising target for DM indirect searches, unless the DM distribution exhibits a large core at its center.
          We introduce a set of representative DM density profiles and compute their $J$-factors.
          We then consider a simple model of thermal relic DM and constrain it with the upper limits on the annihilation cross section obtained by the H.E.S.S. Inner Galaxy Survey~\cite{Abdalla_2022}, assessing the impact of the density profile uncertainty on the constraints.
    
    \item In the third chapter, we study the adiabatic growth of the supermassive black hole at the center of our Galaxy, which redistributes the DM around it into a cusp, called a spike.
          Then we derive the relation between the slope of the spike and that of the DM halo and discuss it for different initial density profiles.
          We introduce two benchmark spike profiles, compute their $J$-factors and finally quantify the impact of the spike on the cross section limits derived in the previous chapter.
     
\end{itemize}


Finally, we summarize and discuss our main results and their implications.

All the code developed and used in this work for numerical calculations is available at \href{https://github.com/IvanDalena/Bachelor-thesis}{\textcolor{githublink}{github.com/IvanDalena/Bachelor-thesis}}.